\section{Evaluation：Measured Capability 取决于 Elicitation}

安全与能力评估看似是给模型出题，实际还包括 prompt、scaffold、tool、time budget、environment、judge 与统计不确定性。\term{elicitation}（能力激发）是通过合理 prompting、tool access、search 或 fine-tuning，让评估尽可能接近模型可达到能力。Elicitation 太弱会低估风险，太强又可能脱离真实 attacker/user 条件。

\subsection{Task、elicitation、environment 与 judge}

评估首先定义 construct：是在测 knowledge、planning、tool use、cyber/CBRN uplift 还是 harmful propensity；随后确定 elicitation budget 和权限；environment 应模拟真实 feedback 和 constraint；judge 则可由专家、人群或模型构成，并需要 rubric、inter-rater agreement 与 calibration。

\lecturefigure{11-evaluation-stack.png}{模型评估的分数由 task、elicitation、environment、judge 和 confidence 共同决定。}{Anthropic evaluation research；本地字幕 27:30--30:40；概念重绘。}

\begin{knowledgebox}{读图：低分可能是能力低，也可能是没激发出来}
固定 prompt 的失败不能证明 model 无能力；允许无限工具和人工搜索的成功也不代表普通用户可复现。Evaluation report 应写 model/version、system prompt、tools、samples、search effort、judge、scoring 和 uncertainty，并区分 observed performance 与 estimated upper bound。
\end{knowledgebox}

\begin{importantbox}{置信区间比单点分数更诚实}
若 $n$ 个独立样本中成功 $k$ 次，经验成功率为
\[
\hat p=\frac{k}{n}.
\]
其中 $\hat p$ 是有限样本估计，真正能力 $p$ 仍有统计区间；当高风险 gate 依赖极低 failure rate 时，几十个样本远远不足。还要处理非独立任务、prompt selection 与 evaluator bias，不能只报告整数百分比。
\end{importantbox}

\begin{warningbox}{高风险评估不应泄露可操作细节}
公开材料可以说明风险类别、评估框架、access control、red-team process 和 aggregate finding，但不应发布能显著降低有害行为门槛的具体流程、材料或绕过方法。Transparency 的目标是可问责，而不是把 evaluation set 变成能力教程。
\end{warningbox}

\subsection{本章小结}

Evaluation 是 measurement system，不是题库。Elicitation、environment 和 judge 决定分数含义；高风险决策需要 uncertainty、external review 和版本化证据。

